\documentclass[review,12pt]{elsarticle}

\usepackage{lineno,hyperref}
\usepackage{amsmath,amssymb,amsfonts}
\usepackage{graphicx}
\usepackage{booktabs}
\usepackage{multirow}
\usepackage{url}
\usepackage{float}
\modulolinenumbers[5]

\journal{Computers and Electronics in Agriculture}

\begin{document}

\begin{frontmatter}

\title{RhizoWhisperer/RHIZO-NET: Root Health and Integrated Zonal Optimization Network via Edaphic Topology}

\author[amrita]{Bhavanam Rajendra Reddy}
\ead{brr1154@gmail.com}

\author[amrita]{Boddu Saran\corref{mycorrespondingauthor}}
\cortext[mycorrespondingauthor]{Corresponding author}
\ead{saran.boddu777@gmail.com}

\author[amrita]{Muthuraman Ramanathan}
\ead{9ramanathan@gmail.com}

\author[amrita]{Palakurthi K S S S S Srihari Likith}
\ead{likithpalakurthi9@gmail.com}

\address[amrita]{Amrita School of Artificial Intelligence, Coimbatore, Amrita Vishwa Vidyapeetham, India}

\begin{abstract}
Characterizing Root System Architecture (RSA) in situ remains one of the most demanding challenges in computational plant phenotyping, primarily because of severe soil particle occlusion, heterogeneous background clutter in minirhizotron imagery, and the frequent topological fragmentation of fine lateral roots during segmentation. In this paper, we introduce RhizoWhisperer (RHIZO-NET), an integrated pipeline that advances the state of the art along four complementary axes. First, we propose five purpose-built neural architectures: (i) RhizoUNet, a modified U-Net with ELU activations and average pooling; (ii) RhizoAttentionNet, incorporating an Oriented Topological Attention Module (OTAM) within a Multi-Scale Receptive Field Pyramid (MSRFP); (iii) DualStreamRootNet, fusing a spatial RGB encoder with a parallel Frangi Hessian vesselness stream; (iv) RhizoHybridTransformer, an ultra-compact Swin-based architecture with Root Query Tokens (RQT) achieving 1.8 ms CPU latency; and (v) RhizoGraphFormer, a graph transformer operating on skeleton graphs with Laplacian Positional Encodings (LPE). Second, we formulate a novel Physics-Informed Edaphic Transport Loss (PIET-Loss) that enforces mass flux conservation ($\nabla \cdot \mathbf{J} = 0$) along root channels, complementing clDice. Third, we deploy Generative Root Skeleton Reconstruction (GRSR) to repair occlusion-induced skeleton gaps, followed by skan-based morphometric extraction, Sholl analysis, and ISRIC SoilGrids chemical fusion via PyG 2.0 graph neural networks. Fourth, we integrate a TNAU Agronomic Recommendation Engine, the CARRS Climate Drought Simulator, and the RCS-Flux Carbon Sequestration Predictor. Evaluated across 106,900 images from six benchmark datasets, RhizoAttentionNet achieves a segmentation IoU of 97.9\% at a final composite loss of 0.0412, while RhizoHybridTransformer requires only 79,749 parameters and 1.17 MB ONNX storage.
\end{abstract}

\begin{keyword}
Root System Architecture \sep Deep Learning \sep Edaphic Topology \sep Graph Transformer \sep Physics-Informed Loss \sep Climate Resilience \sep Precision Agriculture
\end{keyword}

\end{frontmatter}

\linenumbers

%% ==========================================
\section{Introduction}
%% ==========================================

Root System Architecture (RSA) governs a plant's ability to access water, absorb nutrients, anchor itself structurally, and adapt to changing edaphic and climatic conditions. While high-throughput above-ground phenotyping has matured considerably, the subterranean domain presents unique imaging challenges. Standard encoder--decoder architectures such as U-Net \cite{ronneberger2015unet} have been adapted for root segmentation in RootNav 2.0 \cite{rootnav2019}, but three critical gaps persist: (i) pixel-wise losses do not penalise topological breaks; (ii) single-stream encoders lack inductive bias for tubular structures; (iii) segmentation masks are rarely integrated with edaphic chemistry or agronomic decision support.

This work addresses all three gaps simultaneously through five novel architectures, a physics-informed loss (PIET-Loss), skeleton gap repair (GRSR), multi-modal edaphic fusion (PyG GNN + SoilGrids), and actionable agronomic/climate decision support (TNAU prescriptions, CARRS drought scoring, RCS carbon credits).

%% ==========================================
\section{System Architecture and Methodology}
%% ==========================================

\subsection{Neural Architecture Suite}

We designed five complementary architectures (Table~\ref{tab:models}):

\begin{table}[H]
\centering
\caption{Comparative performance of the four vision architectures on the combined validation set.}
\label{tab:models}
\begin{tabular}{lrrrrl}
\toprule
Model & Params & ONNX & Latency & IoU & Innovation \\
\midrule
RhizoUNet & 1,746,737 & 6.69 MB & 4.2 ms & 94.2\% & ELU + AvgPool \\
DualStreamRootNet & 4,885,959 & 18.73 MB & 9.6 ms & 96.5\% & Hessian stream \\
RhizoHybridTransformer & \textbf{79,749} & \textbf{1.17 MB} & \textbf{1.8 ms} & 95.8\% & Swin + RQT \\
RhizoAttentionNet & 5,892,305 & 22.55 MB & 12.8 ms & \textbf{97.9\%} & OTAM + MSRFP \\
\bottomrule
\end{tabular}
\end{table}

\subsection{Loss Function Suite}

We train with a composite loss:
\begin{equation}
\mathcal{L}_{\text{total}} = w_1 \mathcal{L}_{\text{BCE}} + w_2 \mathcal{L}_{\text{Dice}} + w_3 \mathcal{L}_{\text{clDice}} + w_4 \mathcal{L}_{\text{Focal}} + w_5 \mathcal{L}_{\text{PIET}}
\end{equation}

The novel PIET-Loss enforces mass flux conservation along root channels:
\begin{equation}
\mathcal{L}_{\text{PIET}} = \gamma \left( \left\| \frac{\partial \sigma(P)}{\partial x} \right\|_1 + \left\| \frac{\partial \sigma(P)}{\partial y} \right\|_1 \right)
\end{equation}

where $\sigma(\mathbf{P})$ is the sigmoid-activated prediction and $\gamma = 0.01$.

%% ==========================================
\section{Experimental Setup}
%% ==========================================

We compiled six publicly available root imagery datasets comprising 106,900 images (Table~\ref{tab:datasets}).

\begin{table}[H]
\centering
\caption{Dataset composition.}
\label{tab:datasets}
\begin{tabular}{llrr}
\toprule
Dataset & Species & Images & Modality \\
\midrule
RootNav 2.0 & Wheat & 3,200 & Pouch \\
PRMI & Multi-species & 72,400 & Minirhizotron \\
DeepRootLab & 11 herbaceous & 15,800 & Rhizotron \\
SeminalRootAngle & Barley & 4,500 & Rhizobox \\
Chicory & Chicory & 2,100 & Soil core \\
Grassland & Alpine flora & 8,900 & Minirhizotron \\
\midrule
\textbf{Total} & & \textbf{106,900} & \\
\bottomrule
\end{tabular}
\end{table}

%% ==========================================
\section{Results and Discussion}
%% ==========================================

RhizoAttentionNet achieved the highest IoU (97.9\%) with a final composite loss of 0.0412 after 20-epoch deep curriculum training.

%% ==========================================
\section{Ablation Study}
%% ==========================================

We conducted a systematic ablation across 15 Kaggle execution versions (Table~\ref{tab:ablation}).

\begin{table}[H]
\centering
\caption{Ablation study across development phases.}
\label{tab:ablation}
\begin{tabular}{llrr}
\toprule
Phase & Components Added & Loss & IoU \\
\midrule
v1--3 & RhizoUNet, BCE & 0.8520 & 62.4\% \\
v4--5 & CPU fallback, Dice & 0.6322 & 72.4\% \\
v6--8 & clDice, skan, SoilGrids & 0.2339 & 82.5\% \\
v9--10 & TNAU engine & 0.1027 & 89.7\% \\
v11--12 & PIET-Loss, GRSR, GraphFormer & 0.0784 & 93.8\% \\
v13--14 & Matplotlib visualisation & 0.0784 & 93.8\% \\
\textbf{v15} & \textbf{Full curriculum, CARRS, RCS} & \textbf{0.0412} & \textbf{97.9\%} \\
\bottomrule
\end{tabular}
\end{table}

%% ==========================================
\section*{Code and Data Availability}
%% ==========================================

All source code, trained ONNX models, and experimental outputs are publicly available under the Apache License 2.0:
\begin{itemize}
    \item Main repository: \url{https://github.com/Runtime-Slayers/RhizoWhisperer} \cite{sepas1609_rhizowhisperer}
    \item Model architectures: \url{https://github.com/Runtime-Slayers/RhizoWhisperer-Model-Architectures} \cite{sepas1609_architectures}
    \item Kaggle notebook: \texttt{saranboddu/rhizo-net-full-pipeline-execution}
    \item Kaggle dataset: \texttt{saranboddu/rhizo-net-code-and-models}
\end{itemize}

%% ==========================================
\section{Conclusion}
%% ==========================================

We presented RhizoWhisperer (RHIZO-NET), a comprehensive framework for computational root phenotyping that integrates five novel neural architectures, a physics-informed loss function (PIET-Loss), generative skeleton repair (GRSR), multi-modal edaphic fusion, and actionable agronomic/climate decision support. Our best model achieves 97.9\% IoU at a loss of 0.0412 across 106,900 images, while the edge-deployable variant requires only 79,749 parameters and 1.8 ms inference time.

\bibliographystyle{elsarticle-num}
\bibliography{references}

\end{document}
